\documentclass[11pt,a4paper]{article}

% =========================
% Кодировка и язык
% =========================
\usepackage[T2A]{fontenc}
\usepackage[utf8]{inputenc}
\usepackage[russian,english]{babel}

% =========================
% Разметка страницы
% =========================
\usepackage[a4paper,margin=2.5cm]{geometry}
\usepackage{setspace}
\onehalfspacing

% =========================
% Шрифты и математика
% =========================
\usepackage{lmodern}
\usepackage{microtype}
\usepackage{amsmath,amssymb,amsfonts}
\usepackage{mathtools}
\usepackage{bm}

% =========================
% Рисунки и таблицы
% =========================
\usepackage{graphicx}
\usepackage{float}
\usepackage{booktabs}
\usepackage{array}
\usepackage{multirow}
\usepackage{tabularx}
\usepackage{longtable}
\usepackage{caption}
\usepackage{subcaption}

% =========================
% Ссылки и перекрёстные ссылки
% =========================
\usepackage[hidelinks]{hyperref}
\usepackage[nameinlink,noabbrev]{cleveref}

% =========================
% Алгоритмы и код
% =========================
\usepackage{algorithm}
\usepackage{algpseudocode}
\usepackage{listings}
\usepackage{xcolor}

% =========================
% Колонтитулы
% =========================
\usepackage{fancyhdr}
\pagestyle{fancy}
\fancyhf{}
\fancyhead[L]{GRA-Unified}
\fancyhead[R]{\thepage}

% =========================
% Библиография
% =========================
\usepackage[numbers,sort&compress]{natbib}
\bibliographystyle{unsrtnat}

% =========================
% Собственные команды
% =========================
\newcommand{\GRA}{\textsc{GRA}}
\newcommand{\Unified}{\textsc{Unified}}
\newcommand{\code}[1]{\texttt{#1}}
\newcommand{\vect}[1]{\bm{#1}}

% =========================
% Оформление листингов
% =========================
\lstdefinestyle{mystyle}{
    basicstyle=\ttfamily\small,
    keywordstyle=\color{blue},
    commentstyle=\color{gray},
    stringstyle=\color{teal},
    numbers=left,
    numberstyle=\tiny,
    stepnumber=1,
    numbersep=8pt,
    showstringspaces=false,
    breaklines=true,
    frame=single
}
\lstset{style=mystyle}

% =========================
% Название
% =========================
\title{GRA-Unified: семантический фреймворк устойчивости для модульных AI-архитектур}
\author{Ваше имя \and Имя соавтора}
\date{\today}

\begin{document}

\maketitle

\begin{abstract}
В статье представлен GRA-Unified — модульный фреймворк для семантического искусственного интеллекта, предназначенный для иерархического рассуждения, управления устойчивостью и композиции агентных систем. Архитектура объединяет слои представления, контроля и оценки в единую систему для экспериментального исследования продвинутых AI-моделей. Описаны цели проектирования, внутренние модули, стратегия реализации и протокол предварительной оценки. Данный текст служит технической основой для дальнейшего теоретического и экспериментального развития.
\end{abstract}

\noindent\textbf{Ключевые слова:} семантический ИИ, модульная архитектура, координация агентов, управление устойчивостью, иерархическое рассуждение, унифицированный фреймворк

\tableofcontents
\newpage

\section{Введение}
Цель GRA-Unified состоит в создании общей архитектуры для экспериментов с семантическим ИИ, где рассуждение, память, контроль и взаимодействие агентов рассматриваются как отдельные, но совместимые уровни. Эта статья документирует текущую структуру, терминологию и каркас реализации фреймворка. Структура намеренно сделана расширяемой, чтобы в будущих версиях можно было добавлять наборы тестов, более сильные механизмы вывода и более явное отображение теории в код.

Данная работа не претендует на окончательную полноту. Напротив, это воспроизводимое техническое описание системы, находящейся в активной разработке.

\section{Предпосылки}
Современные AI-системы часто объединяют несколько подсистем: модели векторных представлений, retrieval-слой, планировщик, инструменты и мониторинг. В таких системах связь между компонентами становится трудно анализировать, особенно если цель состоит не только в повышении качества задач, но и в обеспечении семантической согласованности и структурной устойчивости. GRA-Unified отвечает на этот вызов через явную организацию системы вокруг интерфейсов и концепций верхнего уровня управления.

\section{Обзор фреймворка}
Фреймворк организован по следующим концептуальным слоям:

\begin{itemize}
    \item \textbf{Семантический слой}: кодирует значения, символы и отношения.
    \item \textbf{Слой управления}: управляет устойчивостью, гейтингом и переходами состояний.
    \item \textbf{Агентный слой}: поддерживает один или несколько автономных единиц рассуждения.
    \item \textbf{Слой памяти}: хранит долговременное состояние, контекст и трассы.
    \item \textbf{Слой оценки}: измеряет согласованность, устойчивость и полезность решений.
\end{itemize}

На рисунке~\ref{fig:architecture} показана упрощённая архитектура системы.

\begin{figure}[H]
    \centering
    \fbox{
    \begin{minipage}[c][5cm][c]{0.88\linewidth}
    \centering
    \vspace{1cm}
    \textbf{Здесь вставляется рисунок архитектуры}\\
    \vspace{0.5cm}
    Семантический слой $\rightarrow$ Слой управления $\rightarrow$ Агентный слой $\rightarrow$ Слой оценки
    \vspace{1cm}
    \end{minipage}
    }
    \caption{Концептуальная архитектура GRA-Unified.}
    \label{fig:architecture}
\end{figure}

\section{Проектирование системы}
Система построена вокруг явных переходов состояний и модульных контрактов. Каждый модуль имеет входы, выходы и внутренние параметры, которые можно отслеживать отдельно. Это облегчает проверку свойств устойчивости и сравнение различных семантических политик в одном и том же вычислительном окружении.

\subsection{Основные модули}
Основные модули включают:

\begin{enumerate}
    \item \textbf{Парсер}: преобразует сырой ввод в структурированное представление.
    \item \textbf{Построитель семантического графа}: создаёт реляционные структуры.
    \item \textbf{Контроллер устойчивости}: регулирует пороги и переходы состояний.
    \item \textbf{Координатор}: распределяет задачи между агентами и подмодулями.
    \item \textbf{Логгер}: записывает состояние, метрики и события.
\end{enumerate}

\subsection{Представление состояния}
Пусть глобальное состояние системы обозначается как \( S_t \). Тогда в момент времени \( t \) система может быть записана как
\[
S_t = \{m_t, c_t, a_t, r_t\},
\]
где \( m_t \) — состояние памяти, \( c_t \) — состояние управления, \( a_t \) — состояние агента, а \( r_t \) — текущая трасса рассуждения.

\section{Скетч реализации}
Реализацию удобно организовать как Python-пакет с отдельными модулями для семантики, управления, агентов, памяти и оценки. Минимальная структура репозитория может выглядеть так:

\begin{lstlisting}[language=bash,caption={Пример структуры репозитория}]
GRA-Unified/
├── paper.tex
├── references.bib
├── src/
│   ├── gra/
│   │   ├── __init__.py
│   │   ├── semantic.py
│   │   ├── control.py
│   │   ├── agent.py
│   │   ├── memory.py
│   │   └── evaluation.py
├── figures/
│   ├── architecture.png
│   └── results.png
└── experiments/
    └── benchmark_runs.py
\end{lstlisting}

\section{Протокол экспериментов}
Для оценки фреймворка рекомендуется следующий протокол:

\begin{itemize}
    \item Сравнивать базовую конфигурацию и унифицированную конфигурацию на одинаковых задачах.
    \item Измерять согласованность, стабильность ответа и качество выполнения задачи.
    \item Публиковать средние значения, дисперсии и типовые ошибки.
    \item Проводить абляционные тесты для семантических и управляющих модулей.
\end{itemize}

Таблица~\ref{tab:metrics} показывает пример шаблона отчёта о результатах.

\begin{table}[H]
    \centering
    \caption{Пример метрик оценки для GRA-Unified.}
    \label{tab:metrics}
    \begin{tabular}{lccc}
        \toprule
        Конфигурация & Согласованность $\uparrow$ & Устойчивость $\uparrow$ & Балл задачи $\uparrow$ \\
        \midrule
        Базовая & 0.71 & 0.64 & 0.68 \\
        Unified без контроля & 0.75 & 0.59 & 0.70 \\
        GRA-Unified полная & 0.83 & 0.81 & 0.79 \\
        \bottomrule
    \end{tabular}
\end{table}

\section{Результаты}
Раздел результатов должен содержать эмпирические выводы, визуализированные при помощи графиков и таблиц. Если доступны реальные данные бенчмарков, сюда можно вставить графики сходимости, абляционные диаграммы или траектории устойчивости.

\begin{figure}[H]
    \centering
    \fbox{
    \begin{minipage}[c][5cm][c]{0.88\linewidth}
    \centering
    \vspace{1cm}
    \textbf{Здесь вставляется график результатов}\\
    \vspace{0.5cm}
    Например: устойчивость во времени, сравнение бенчмарков или абляционный график
    \vspace{1cm}
    \end{minipage}
    }
    \caption{Пример заглушки для рисунка с результатами.}
    \label{fig:results}
\end{figure}

\section{Обсуждение}
Главное преимущество фреймворка состоит в том, что он отделяет семантическую организацию от низкоуровневого выполнения. Это повышает интерпретируемость и позволяет понять, на каком этапе возникает ошибка. Второе преимущество — расширяемость: новые модули рассуждения можно добавлять без переписывания всей системы.

Ограничение состоит в том, что фреймворку всё ещё нужны формально определённые метрики устойчивости и наборы тестовых задач. Без этого система остаётся сильным архитектурным предложением, но не полностью верифицированной экспериментальной платформой.

\section{Заключение}
GRA-Unified предлагает структурированную основу для построения семантических AI-систем с явным контролем устойчивости и модульной композицией. Шаблон статьи подходит для описания репозитория, отчёта об экспериментах и связи теории с реализацией. В будущих работах следует добавить формальные метрики, воспроизводимые эксперименты и более глубокие сравнения с альтернативными архитектурами.

\section*{Благодарности}
Здесь можно разместить благодарности.

\bibliography{references}

\end{document}